\documentclass{miniopenclaw-poster}

% Optional quick theme adjustments:
% \PosterSetupColors{174A7C}{EAF2F8}{F5F5F5}
% \renewcommand{\PosterTitleFont}{\Huge\bfseries}
% \renewcommand{\PosterBlockTitleFont}{\Large\bfseries}
\renewcommand{\PosterBlockInnerYSep}{0.35cm}
\renewcommand{\PosterBlockTopBottomSpace}{0.16cm}
\graphicspath{{uploads/}}

\pgfplotsset{
  mockupchart/.style={
    width=0.285\linewidth,
    height=4.8cm,
    ymajorgrids=true,
    grid style={white},
    axis background/.style={fill=postergray},
    tick label style={font=\scriptsize},
    label style={font=\scriptsize},
    title style={font=\small\bfseries, align=center},
    xticklabel style={font=\scriptsize},
    ymin=0,
    enlarge x limits=0.25,
  }
}

\newcommand{\ArchitectureFigure}{%
  \begin{center}
    \begin{tikzpicture}[>=stealth', thick, scale=0.9, transform shape]
      \tikzset{
        bnode/.style={rectangle, draw=posterblue, fill=blue!10, text width=9.2em, align=center, rounded corners, minimum height=3.3em, font=\small},
        smallnode/.style={rectangle, draw=posterblue, fill=postergray, text width=8.5em, align=center, rounded corners, minimum height=3em, font=\scriptsize},
        endpoint/.style={font=\small\bfseries, align=center},
        line/.style={draw=posterblue, -latex', shorten >=2pt, shorten <=2pt},
        side/.style={draw=posterblue, -latex', shorten >=2pt, shorten <=2pt, rounded corners=6pt}
      }

      % Nodes
      \node[endpoint] (user) at (0,0) {User\\Message};
      \node[bnode] (orchestrator) at (0,-1.25) {\textbf{Orchestrator}\\controls the run};
      \node[bnode] (planner) at (0,-2.85) {Planner\\ReAct step};
      \node[bnode] (policy) at (0,-4.45) {Policy Engine\\classifies action};
      \node[bnode] (executor) at (0,-6.05) {Executor\\validates \& runs tools};
      \node[smallnode] (memory) at (-4.3,-3.65) {Memory\\retrieval \& storage};
      \node[smallnode] (audit) at (-4.3,-5.25) {Audit Log\\append-only trace};
      \node[smallnode] (events) at (4.3,-4.45) {SSE Events\\frontend updates};
      \node[endpoint] (answer) at (0,-7.55) {Final Answer\\or Next Iteration};

      % Edges
      \path [line] (user) -- (orchestrator);
      \path [line] (orchestrator) -- (planner);
      \path [line] (planner) -- (policy);
      \path [line] (policy) -- (executor);
      \path [side] (executor.west) -- ++(-1.2,0) |- node[pos=0.28, left, font=\scriptsize] {observe} (orchestrator.west);
      \path [side] (orchestrator.west) -- ++(-1.2,0) |- (memory.north);
      \path [line] (memory) -- (audit);
      \path [side] (orchestrator.east) -- ++(1.2,0) |- (events.north);
      \path [side] (executor.east) -- ++(1.2,0) |- (events.south);
      \path [line] (executor) -- (answer);
      \path [side] (events.south) |- (answer.east);
    \end{tikzpicture}
  \end{center}
}
\newcommand{\EvaluationGraphs}{%
  \begin{center}
  \begin{tabular}{@{}c@{\hspace{0.2cm}}c@{\hspace{0.2cm}}c@{}}
    \begin{tikzpicture}
      \begin{axis}[
        mockupchart,
        title={Erfolgsrate},
        ylabel={Success (\%)},
        symbolic x coords={Baseline,Plan+ReAct,Full},
        xtick=data,
        xticklabel style={rotate=25, anchor=east},
        ymax=100,
        bar width=12pt,
      ]
        \addplot+[ybar, fill=posterblue!70, draw=posterblue] coordinates {(Baseline,58) (Plan+ReAct,73) (Full,86)};
      \end{axis}
    \end{tikzpicture}
    &
    \begin{tikzpicture}
      \begin{axis}[
        mockupchart,
        title={Avg. Tool Calls},
        ylabel={Calls / run},
        symbolic x coords={Baseline,Plan+ReAct,Full},
        xtick=data,
        xticklabel style={rotate=25, anchor=east},
        ymax=10,
        bar width=12pt,
      ]
        \addplot+[ybar, fill=blue!35, draw=posterblue] coordinates {(Baseline,6.9) (Plan+ReAct,5.4) (Full,4.2)};
      \end{axis}
    \end{tikzpicture}
    &
    \begin{tikzpicture}
      \begin{axis}[
        mockupchart,
        title={Cost vs. Success},
        xlabel={Cost / run (\$)},
        ylabel={Success (\%)},
        xmin=0.00, xmax=0.10,
        ymin=45, ymax=95,
        xtick={0.02,0.04,0.06,0.08},
      ]
        \addplot+[only marks, mark=*, mark size=3pt, color=posterblue] coordinates {(0.025,58) (0.043,73) (0.068,86)};
        \node[font=\scriptsize, anchor=south west] at (axis cs:0.025,58) {Baseline};
        \node[font=\scriptsize, anchor=south west] at (axis cs:0.043,73) {Plan+ReAct};
        \node[font=\scriptsize, anchor=south west] at (axis cs:0.068,86) {Full};
      \end{axis}
    \end{tikzpicture}
  \end{tabular}
  \end{center}
}

\begin{document}

% ============================================================
% A3-Blatt 1: Projektüberblick
% ============================================================
\PosterTitle
  {Mini-OpenClaw {\mdseries\RobotEmoji}}
  {Held Johannes (11705340 / e11705340@student.tuwien.ac.at) \quad | \quad Ibrahim Ahmad (11826186 / e11826186@student.tuwien.ac.at)}
  {Ibrar Muhammad ([STUDENT-ID] / [EMAIL\_ADDRESS]) \quad | \quad Bourbigua Manal ([STUDENT-ID] / [EMAIL\_ADDRESS])}
  {Gruppe 09 \quad | \quad Version 0.1.0 \quad | \quad 194.211 Applied Generative AI and LLM-based Systems, TU Wien \quad}

\PosterBlock{\SectionIcon{info-circle} Project Overview}{
  Mini-OpenClaw is a lightweight, local-first AI agent that converts natural-language requests into safe, auditable tool executions on a local machine. The project emphasizes transparent, inspectable pipelines for agent-based interactions and is developed as part of the Applied Generative AI course at TU Wien.
}

\begin{multicols}{2}

\ArchitectureFigure

\PosterBlock{\SectionIcon{lightbulb} Core Idea}{
  \textbf{Hybrid Plan $\rightarrow$ ReAct $\rightarrow$ Replan:} The LLM first creates a goal checklist, then executes actions in an iterative reason--act--observe loop, and replans when progress stalls or goals are skipped.
}

\PosterBlock{\SectionIcon{bullseye} Goals}{
  \begin{itemize}
    \item Make local AI-agent actions safe and auditable.
    \item Separate planning, policy checks, execution, memory, and logging.
    \item Provide real-time UI feedback and approval gates for risky steps.
  \end{itemize}
}
\PosterBlock{\SectionIcon{tools} Tools / Skills}{
  Tools are registered through a SkillRegistry: file operations, safe shell execution, memory search/storage, sub-agent delegation, guarded web fetches, task scheduling, and run explanations.
}
\end{multicols}

\begin{multicols}{2}
\PosterBlock{\SectionIcon{server} Backend: \texttt{apps/api/}}{
  \begin{itemize}
    \item FastAPI service for agent operations and API routes.
    \item Core components: Orchestrator, Planner, Policy Engine, Executor, Audit Logger, Memory Manager, and Task Scheduler.
    \item Pluggable LLM provider abstraction for Anthropic, Google Gemini, and Ollama.
    \item SQLite persistence for runs, memory items, and scheduled tasks.
  \end{itemize}
}

\PosterBlock{\SectionIcon{desktop} Frontend: \texttt{apps/web/}}{
  \begin{itemize}
    \item React/TypeScript interface for interacting with the agent.
    \item Real-time Server-Sent Events for run status, plans, approvals, and traces.
    \item Main UI components: ChatPanel, PlanPreview, ExecutionGraph, ApprovalCard, ToolTrace, RunHistory, MemoryBrowser, and SchedulerPage.
  \end{itemize}
}
\end{multicols}

\newpage

% ============================================================
% A3-Blatt 2: Architektur und Ausführung
% ============================================================
\PosterPageTitle{Architecture, Runtime and Mockup Evaluation}

\begin{multicols}{2}
\PosterBlock{\SectionIcon{project-diagram} Execution Flow}{
  \begin{enumerate}
    \item User message starts an Orchestrator-managed run.
    \item Planner selects actions in a reason--act--observe loop.
    \item Policy Engine classifies actions as safe, approval-required, or forbidden.
    \item Executor validates tools, retries transient errors, and returns observations.
    \item Memory, Audit Logger, and SSE events persist context and stream progress.
  \end{enumerate}
}

\PosterBlock{\SectionIcon{brain} Memory System}{
  Mini-OpenClaw uses five memory layers: \textbf{facts}, \textbf{episodes}, \textbf{summaries}, user-reviewed \textbf{strategies}, and \textbf{preferences}. After every \texttt{DREAM\_INTERVAL} runs, Agent Dreams proposes workflow insights as \texttt{pending\_review}; users can accept, dismiss, or edit them.

  Retrieval combines \textbf{70\% vector similarity} with local \texttt{all-MiniLM-L6-v2} embeddings and \textbf{30\% SQL keyword matching}; relevant context is injected before each planning call.
}
\columnbreak

\PosterBlock{\SectionIcon{star} Key Features}{
  \begin{itemize}
    \item Sub-agent delegation and confidence-gated clarification.
    \item Scheduled one-time or recurring tasks with approval options.
    \item Real-time SSE streaming, status updates, and execution graphs.
    \item Budget-aware replanning, retries, and self-reflection quality gates.
  \end{itemize}
}

\PosterBlock{\SectionIcon{shield-alt} Security Model}{
  \begin{itemize}
    \item Only registered tools with valid JSON schemas can be invoked.
    \item Policy engine enforces workspace boundaries, shell allowlists, and injection detection.
    \item Risky writes, shell commands, web fetches, and delegation require approval.
    \item Append-only audit logging records decisions and actions.
  \end{itemize}
}

\end{multicols}

\PosterBlock{\SectionIcon{chart-line} Mockup Evaluation}{
  \EvaluationGraphs
  \textbf{Illustrative benchmark design:} compare a baseline agent, the Plan+ReAct loop, and the full Mini-OpenClaw pipeline on representative local automation tasks. The mockup metrics show how explicit planning can improve \textbf{Erfolgsrate}, reduce average tool-call overhead, and make the cost--quality trade-off visible.
}

\newpage

% ============================================================
% A3-Blatt 3: Features, Sicherheit und Ausblick
% ============================================================
\PosterPageTitle{Features, Security and Outlook}

\begin{multicols}{2}
\PosterBlock{\SectionIcon{flask} Evaluation Plan}{
  \begin{itemize}
    \item Build a small task suite for file edits, shell reads, memory lookup, and scheduling.
    \item Report success rate, mean tool calls, approval frequency, latency, and estimated cost.
    \item Use failed runs to refine policies, prompt templates, and replanning triggers.
  \end{itemize}
}

\PosterBlock{\SectionIcon{forward} Current Status and Next Steps}{
  \begin{itemize}
    \item Project status: ongoing, version 0.1.0.
    \item Continue evaluating reliability of replanning, memory retrieval, and tool approvals.
    \item Improve frontend explanations and refine threat-model documentation in \texttt{docs/threat-model.md}.
  \end{itemize}
}
\end{multicols}

\PosterBlock{\SectionIcon{check-circle} Contribution}{
  \textbf{Mini-OpenClaw demonstrates a practical design for transparent local agents:} planning is explicit, execution is policy-gated, memory is inspectable, and every important decision is auditable. This makes the system suitable for experimenting with LLM-based automation while preserving user oversight.
}


\PosterBlock{\SectionIcon{image} UI Mockup Screenshot}{
  \begin{center}
    \includegraphics[width=0.82\linewidth,height=10.2cm,keepaspectratio]{Bildschirmfoto 2026-06-01 um 14.50.06.png}
  \end{center}
}


\end{document}
